% Generated draft: requires mathematical and source review.
\documentclass[11pt]{article}
\usepackage[margin=25mm]{geometry}
\usepackage{fontspec}
\setmainfont{Latin Modern Roman}
\usepackage{amsmath,amssymb,mathtools,unicode-math}
\setmathfont{Latin Modern Math}
\usepackage{enumitem,needspace,xurl,hyperref}
\setlength{\parindent}{0pt}
\setlength{\parskip}{5pt}
\setlength{\emergencystretch}{3em}
\newcommand{\N}{\mathbb{N}}
\newcommand{\Z}{\mathbb{Z}}
\newcommand{\Q}{\mathbb{Q}}
\newcommand{\R}{\mathbb{R}}
\newcommand{\C}{\mathbb{C}}
\newcommand{\F}{\mathbb{F}}
\newcommand{\A}{\mathbb{A}}
\newcommand{\PP}{\mathbb{P}}
\newcommand{\E}{\mathbb{E}}
\newcommand{\eps}{\varepsilon}
\newcommand{\dd}{\,\mathrm{d}}
\begin{document}
\section*{UnsolvedMath ID: 30002540}
\textbf{Source problem code:} catalogue entry
\subsection{Definitions and mathematical statement}
The supplied problem is a research programme concerning Feynman categories, Tannakian reconstruction, Hopf cyclic cohomology, and Hopf algebroids. It does not specify enough data to determine a single quantified existence assertion.

For terminology only, let \(k\) be a field. In the neutral convention, a Tannakian category over \(k\) is a rigid abelian symmetric monoidal \(k\)-linear category \(\mathcal T\), with \(\operatorname{End}_{\mathcal T}(\mathbf 1)\cong k\), admitting a faithful exact \(k\)-linear symmetric monoidal functor
\[
\omega:\mathcal T\longrightarrow\operatorname{Vect}^{\mathrm{fd}}_k.
\]
Such a functor is called a fiber functor; \(\operatorname{Vect}^{\mathrm{fd}}_k\) denotes the category of finite-dimensional \(k\)-vector spaces. This convention explains the terminology but is not an additional hypothesis of the source: the record does not specify a field, neutrality, or the target of its fiber functors.

The source asks for a construction or conceptual description addressing the following linked questions:
\begin{enumerate}
\item How can fiber functors and their associated Tannakian categories be exhibited directly in the Feynman-category picture?
\item What role does Hopf cyclic cohomology play in this relationship?
\item Is there a categorification of this relationship connecting it to Hopf algebroids?
\end{enumerate}
Here the relevant Feynman categories, the meaning of a direct construction, the Hopf cyclic complexes and their coefficients, and the intended categorification are unspecified. Likewise, the base algebras and convention for Hopf algebroids are not supplied. No universal quantifier over Feynman categories, equivalence of categories, reconstruction theorem, or cohomological comparison is asserted by this draft. Specifying those data and the required compatibility properties is necessary before these questions can become a precise theorem-shaped target.
\subsection{Short English statement}
The problem asks how to see fiber functors and Tannakian categories directly through Feynman categories, how Hopf cyclic cohomology participates, and whether a categorification links the picture to Hopf algebroids. The supplied record leaves the input data and required constructions unspecified, so this remains a research programme requiring formulation and mathematical review.
\subsection{Sources}
\begin{itemize}
\item Supplied local registry record: \texttt{data/unsolvedmath/problems.json}, permanent UnsolvedMath ID \(30002540\), source problem code \texttt{catalogue entry}, title ``Tannakian Structures in Feynman Categories.'' The recorded proposed year is 2014; no individual proposer is identified.
\item Report attribution: \emph{Mini-Workshop: Batalin-Vilkovisky Algebras, Operads, and Hopf Algebroids} (2014), Oberwolfach Reports, pp.\ 595--634, DOI \texttt{10.4171/owr/2014/11}. Supplied link: \url{https://doi.org/10.4171/owr/2014/11}. The record describes these reports as documenting mathematical workshops and their open-problem discussions; no exact problem page is supplied.
\item Supplied literature reference: Kaufmann and Ward, 2014, \url{https://arxiv.org/abs/1312.1269}.
\item Supplied literature reference: Kowalzig and Kr\"ahmer, 2018, \url{https://arxiv.org/abs/1803.09194}.
\item Supplied literature reference: B\"ohm, 2010, \url{https://arxiv.org/abs/0907.1578}.
\item Literature assessment checked at \texttt{2026-08-21}: the supplied curation labels the problem open and states that Tannakian, Hopf-cyclic, and Hopf-algebroid frameworks exist separately, while no general direct Feynman-category construction answering all three questions is known. This is supplied dated curation derived from the OWR record, not an independently verified literature assessment. The record expressly requests independent verification of time-sensitive status and cited claims.
\end{itemize}
\end{document}
